Beide Monde umkreisen Jupiter. Nach dem dritten Kepler'schen Gesetz ist $T^2/a^3$ für beide gleich:
\begin{align}
 \frac{T_\text{E}^2}{a_\text{E}^3} &= \frac{T_\text{I}^2}{a_\text{I}^3}
\end{align}
Damit ist die Umlaufzeit von Europa
\begin{align}
 T_\text{E} &= \TEuF \\
   &= \TIo\cdot\left(\frac{\aEuO}{\aIoO}\right)^{3/2} \\
   &= \TEu \approx \result{\TEuP{0}{3}}
\end{align}
Europa braucht genau doppelt so lange wie Io (Bahnresonanz).
